\documentclass[11pt, oneside]{article} 
\usepackage{geometry}
\geometry{letterpaper} 
\usepackage{graphicx}
	
\usepackage{amssymb}
\usepackage{amsmath}
\usepackage{parskip}
\usepackage{color}
\usepackage{hyperref}

\graphicspath{{/Users/telliott/Dropbox/Github-Math/figures/}}
% \begin{center} \includegraphics [scale=0.4] {gauss3.png} \end{center}

\title{Angle bisector in a right triangle}
\date{}

\begin{document}
\maketitle
\Large

%[my-super-duper-separator]

\subsection*{stacked triangles}

\label{sec:stacked_triangles}

\begin{center} \includegraphics [scale=0.5] {angle_bisector_r8.png} \end{center}

We stack two right triangles so that the hypotenuse of the first is the base of the second..  We draw the altitude to form another right angle.  

We immediately see an angle equality as marked by the black dots, most simply obtained by using vertical angles to see that we have similar right triangles.  Alternatively, compute sums, adding up the angles in the upper triangle compared to the one whose side is the altitude.

We draw a box and work our way through the diagram marking the other angles.

\begin{center} \includegraphics [scale=0.5] {angle_bisector_r9.png} \end{center}
 
This figure will become important later in the derivation of what are called the sum of angles formulas.  The relationships can be proven using the alternate interior angles theorem.

\begin{center} \includegraphics [scale=0.5] {angle_bisector_r10.png} \end{center}

One more point.  At the beginning of the book is a proof of Thales' theorem about right angles in a semicircle.  We will go through this more carefully later (\hyperref[sec:Thales_theorem]{\textbf{here}}), but assume it for the moment.

Draw the circle with $PS$ as its diameter, then by the (converse of) the theorem, points $R$ and $Q$ are on the circle, because they form right angles.  So $PR$ and $QS$ are secants of the circle.

Looking at the right-hand side, we have labeled the sides and hypotenuse of two similar triangles.  By the properties of similar triangles we have that
\[ \frac{a}{h} = \frac{A}{H} \]
\[ aH = Ah \]

This is a proof that when two secants cross in a circle, the parts of the secants multiply to give two equal products.

\subsection*{angle bisector}

\label{sec:angle_bisector}

We now consider a classic problem:  the ratio of sides when there is an angle bisector.

\begin{center} \includegraphics [scale=0.4] {angle_bisector_r1.png} \end{center}

Suppose the large triangle is a right triangle.  We draw a line joining the vertex on the left with the side opposite. 

This line could in general be drawn anywhere, however two interesting cases are when  the side opposite is bisected (left panel), or when the angle at the left is bisected (right panel).  These two cases are not the same.  In the first $\phi \ne \theta$ and in the second, $c \ne d$.

Let's investigate the second possibility, equal angles.  We are in a position to prove an important theorem.

\subsection*{angle bisector theorem}

\label{sec:angle_bisector_theorem}

With reference to the right-hand figure above, we are to prove that
\[ \frac{d}{b} = \frac{c}{a} \]

$\bullet$ \ The sides and bases are in proportion for a right triangle with bisected angle.

\emph{Proof}.

Draw an altitude for the upper of the two small triangles, meeting the side of length $b$.

\begin{center} \includegraphics [scale=0.4] {angle_bisector_r2.png} \end{center}

The red triangle and the one directly above it are congruent (right panel).  They share a side (the hypotenuse of each), and they are right triangles with the same smaller angle $\theta$.  Therefore, the altitude we just drew has length $c$.

The small triangle with sides $c$ and $d$ (at the top) is similar to the original large triangle.  The reason is that they are both right triangles containing the smaller angle $\gamma$.

By similar triangles, we form equal ratios of the angle opposite $\gamma$ to the hypotenuse:

\[ \frac{a}{b} = \frac{c}{d} \]

This is rearranged simply to give
\[ \frac{d}{b} = \frac{c}{a} \]
which is what we were asked to prove.

The result can be pushed a little further:
\[ \frac{a}{b} = \frac{c}{d} \]
Add $1$ to both sides:
\[ \frac{a + b}{b} = \frac{c + d}{d} \]
\[ \frac{a + b}{c + d} = \frac{b}{d} = \frac{a}{c} \]

$\square$

\begin{center} \includegraphics [scale=0.4] {angle_bisector_r5.png} \end{center}

which is a surprising result and becomes important later in looking at Archimedes method for approximating the value of $\pi$. 

\subsection*{midpoint theorem}

\label{sec:right_triangle_midpoint_theorem}

The line which bisects the hypotenuse in a right triangle is called the median, but this is commonly called the midpoint theorem.  It is a specific case of the theorem for a right triangle.

\begin{center} \includegraphics [scale=0.35] {rt_tri_bisector.png} \end{center}

In a right triangle, draw the line segment from the vertex that contains a right angle to the midpoint of the hypotenuse, separating it into two equal lengths $a$.  We will show that the length of the bisector is also $a$.

\emph{Proof}.

In the right panel, draw the perpendicular from the midpoint $S$ to the base $PR$.  The triangle $SQR$ is similar to the original right triangle by AAA.

Hence the two parts of the base are equal (labeled $b$), because $a/2a = b/2b$.  

Therefore we have two congruent triangles:  $SQR$ and $PQS$ (by SAS).  So the bisector $PS$ is equal in length to $SR$.

Both of the new isosceles triangles formed by the original dashed line have equal base angles.

$\square$

The length $a$ is the radius of the circle, centered at $S$, that contains the three vertices of the original right triangle.

An alternative proof of the midpoint theorem would be to start with any right triangle, and use the hypotenuse as the diameter of a circle.  The center of the circle is on the midpoint of the hypotenuse.

By the converse of Thales circle theorem, the vertex with the right angle is also on the circle.  Therefore, the line connecting that vertex with the midpoint is a radius.

\end{document}